\section{Methodology}

\subsection{Official Metric Contract}

As introduced in Section~1, \gittaskbench emits official \processmetric and
\resultmetric fields for every graded trace. Missing outputs or test-script
failures yield \processmetric$=$False, \resultmetric$=$False. When a trace
reaches grading, the task script writes final \processmetric/\resultmetric
records plus evaluator comments.

The split is meaningful because different task scripts impose qualitatively
different criteria. \texttt{Trafilatura\_01} requires $F1 \ge 0.90$.
\texttt{Trafilatura\_02} requires edit-distance ratio $\le 0.50$.
\texttt{SpeechBrain\_05} requires exact JSON \texttt{prediction} match.
\texttt{PDFPlumber\_02} requires $75\%$ cell-level accuracy. These thresholds
separate reaching grading from meeting the full task criterion, and they vary
enough across families that a single global pass rate blends qualitatively
different evaluation standards.

We interpret \processmetric conservatively throughout the paper.
\processmetric$=$True means that the trace received an official benchmark
record, not that it is a clean semantic near-miss: some official comments
report threshold misses, while others report parsing or evaluation errors after
grading begins. All paper-side aggregates are derived only from these official
outputs; no hand relabeling is involved.

\subsection{Evaluation Cohorts}

We organize the evidence into four cohorts with an intentional hierarchy:

\paragraph{Comparable open-model cohort (primary evidence).}
Eight models evaluated on the full official \gittaskbench task set through the
same agent harness and evaluator contract:
\texttt{Qwen2.5-7B},
\texttt{Qwen2.5-14B},
\texttt{Qwen2.5-32B},
\texttt{Qwen3-14B},
\texttt{Qwen3-32B},
\texttt{DeepSeek-R1-D8B},
\texttt{Mistral-7B}, and
\texttt{Mistral-Small-24B}.

\paragraph{Shared Trafilatura slice (cross-model existence proof).}
Task-aligned runs on \texttt{Trafilatura\_01} and \texttt{Trafilatura\_02} that
extend coverage to \texttt{gpt-5.2} and \texttt{Gemini-2.5-Pro}. This slice
checks whether official success patterns extend beyond the open-model family
but is not comparable in the same sense as the main cohort.

\paragraph{Targeted provider screen (robustness check).}
\texttt{gpt-5-chat}, \texttt{gpt-5.2}, and \texttt{Gemini-2.5-Pro} on a
structured task subset covering the main success-bearing families, the main
\processmetric$=$True/\resultmetric$=$False families, and one image plus one
video stressor. Reported separately from the comparable cohort.

\paragraph{Full-benchmark provider sweeps (extrapolation test).}
The same three provider models on the full benchmark, testing whether targeted
screen activations persist once the full task distribution is restored. Also
reported separately.

\subsection{Aggregation and Reporting}

Each model-task trace is reduced to its official
\processmetric/\resultmetric pair. Family labels come from the task prefix
(e.g., \texttt{Eparse\_02} $\to$ \texttt{Eparse}). Domain labels are fixed at
the family level: \texttt{Eparse}, \texttt{PDFPlumber}, and \texttt{PyPDF2}
are office-document tasks; \texttt{Faker}, \texttt{InvisibleWatermark}, and
\texttt{Stegano} are security tasks; \texttt{Trafilatura} and \texttt{Scrapy}
are web tasks; and so on. The taxonomy is fixed before results, not learned
from them.

This mapping matters because the benchmark is structurally imbalanced. Image
tasks occupy the largest share of evaluation mass, while video is only a small
residue. A global pass rate over this mixture blends very different task
families and very different amounts of evaluation mass.

We use three reporting views throughout the paper:
\begin{enumerate}[leftmargin=1.25em]
\item per-model official totals,
\item domain or family slices aggregated from official outputs,
\item representative official comments from the
\processmetric$=$True, \resultmetric$=$False regime.
\end{enumerate}

The four cohorts are never merged into one aggregate. The comparable cohort
supports the paper's main claims; the provider cohorts are robustness checks
reported in their own sections. All tables and figures are regenerated
deterministically from exported summaries of the official outputs.
